\documentclass[11pt,a4paper]{article}
\usepackage{amsmath,amssymb,amsthm}
\usepackage{geometry}
\usepackage{hyperref}
\usepackage{booktabs}
\usepackage{graphicx}
\usepackage{listings}
\usepackage{xcolor}
\usepackage{longtable}
\usepackage{xspace}

\geometry{margin=1in}

\title{\textbf{A Minigrid-Graph Sudoku Solver}\\
\large Holistic Design, Exact Support Pruning, and Analysis Workflow}
\author{}
\date{}

\lstset{
    basicstyle=\ttfamily\small,
    keywordstyle=\color{blue},
    commentstyle=\color{gray},
    breaklines=true,
    frame=single
}

\newcommand{\optionalfigure}[3]{
    \begin{figure}[h]
        \centering
        \IfFileExists{#1}{
            \includegraphics[width=#2]{#1}
        }{
            \fbox{\parbox[c][2.2in][c]{0.8\textwidth}{\centering Figure placeholder:\\\texttt{\detokenize{#1}}\\Generate with \texttt{scripts/generate\_figures.py}.}}
        }
        \caption{#3}
    \end{figure}
}

\begin{document}

\maketitle

\begin{abstract}
This paper describes the full \texttt{sudoku\_solver} project as an end-to-end Sudoku research and engineering system. The solver decomposes a $9 \times 9$ board into minigrids, generates valid box-level permutations, constructs a compatibility graph across dependent minigrids, prunes unsupported configurations, and extracts complete solutions. In addition to the core algorithm, the project includes dataset analysis tooling, reproducible figure generation, and a paper pipeline intended to evolve as new measurements and scripts are added. The document unifies the previously isolated branch-free relationship note into a broader project view.
\end{abstract}

\section{Project Overview}

The repository is organized around five solver phases and a paper-analysis pipeline:
\begin{enumerate}
    \item Mask initialization
    \item Minigrid permutation generation
    \item Compatibility graph construction
    \item Exact support pruning
    \item Solution extraction
\end{enumerate}

Supporting infrastructure includes sample datasets, CSV analysis outputs, plotting scripts, and LaTeX sources.

\section{Motivation}

Traditional Sudoku solvers often rely on direct backtracking over cells. This project instead treats each $3 \times 3$ minigrid as the primary combinational unit. That shift makes it possible to:
\begin{itemize}
    \item exploit box-local permutation structure,
    \item precompute row and column compatibility masks,
    \item reason about global consistency through a graph of minigrid permutations,
    \item analyze solver behavior phase-by-phase across puzzle datasets.
\end{itemize}

\section{Architecture}

\subsection{Phase 1: Mask Initialization}

The board is parsed into row, column, and box conflict masks. These masks provide the legality checks used during local permutation generation.

\subsection{Phase 2: Minigrid Permutation Generation}

For each minigrid, the solver runs a depth-first search with a minimum-remaining-values heuristic. The output of this phase is a set of valid candidate fillings for each minigrid. These candidates are represented as permutation nodes containing both cell assignments and row/column masks.

\subsection{Phase 3: Graph Construction}

Each permutation node becomes a graph vertex. Edges connect vertices in dependent minigrids when row or column masks are compatible. The project retains a specialized branch-free relationship primitive, documented separately in \texttt{minigrid\_relationship.tex}, for identifying whether two minigrids share row or column constraints.

\subsection{Phase 4: Exact Support Pruning}

Earlier connectivity-only pruning was too aggressive on realistic puzzles because local degree conditions do not imply membership in a globally consistent full-board assignment. The current implementation instead performs exact support pruning: it explores compatible minigrid configurations, marks permutations that appear in at least one full assignment, and discards unsupported nodes. This makes pruning semantically aligned with actual solvability.

\subsection{Phase 5: Solution Extraction}

After pruning, the remaining graph is searched for complete assignments. Valid assignments are reconstructed into Sudoku boards and classified as unsolvable, unique, or ambiguous.

\section{Analysis Pipeline}

The project includes a structured analysis path:
\begin{itemize}
    \item \texttt{examples/analyze\_dataset.rs} generates CSV outputs,
    \item \texttt{results/} stores analysis outputs,
    \item \texttt{scripts/generate\_figures.py} produces paper figures,
    \item \texttt{scripts/run\_paper\_pipeline.sh} ties analysis, plotting, and paper generation together.
\end{itemize}

Current analysis outputs include:
\begin{itemize}
    \item per-difficulty CSV files,
    \item \texttt{complete\_analysis.csv},
    \item \texttt{phase\_breakdown.csv},
    \item \texttt{statistics\_summary.txt}.
\end{itemize}

\section{Generated Figures}

\optionalfigure{figures/permutation_distribution.pdf}{0.85\textwidth}{Distribution of minigrid permutation counts.}

\optionalfigure{figures/graph_reduction.pdf}{0.9\textwidth}{Average graph size before and after pruning.}

\optionalfigure{figures/phase_breakdown.pdf}{0.85\textwidth}{Average time spent in each phase by difficulty.}

\optionalfigure{figures/solution_classification.pdf}{0.8\textwidth}{Solution classification across analyzed puzzles.}

\section{Current Status}

The implementation now includes all five phases and integration tests for representative easy and medium sample puzzles. The paper pipeline is intended to remain open-ended: additional scripts, more detailed analysis, and further figures can be added without restructuring the repository again.

\section{Future Work}

Near-term work includes:
\begin{itemize}
    \item larger dataset runs,
    \item optimization of exact support pruning for harder puzzles,
    \item richer benchmark suites,
    \item expanded plots and tables for publication.
\end{itemize}

Longer-term work may include new data analysis scripts, solver variants, or comparative baselines against conventional cell-level backtracking approaches.

\section{Conclusion}

The repository has evolved from a narrowly focused algorithm note into a complete solver-and-analysis project. The full paper source introduced here is meant to become the main narrative document for that broader effort, while the original minigrid relationship note remains a focused reference on one important low-level optimization.

\end{document}
